\section{Local power closure below a singular cardinal}

\import{cardinal/infinite-exponentiation.tex}
\import{cardinal/regularity.tex}
\import{stationary/club.tex}

\subsection{GCH points and power closure}

\begin{definition}\label{gch_point}
    $\gchPoint{\mu}$ iff $\mu$ is an infinite cardinal and
        $\cardPower{\cardTwo}{\mu}=\cardSucc{\mu}$.
\end{definition}

\begin{definition}\label{gch_points}
    $\gchPoints{\kappa}
        =\{\mu\in\kappa\mid \gchPoint{\mu}\}$.
\end{definition}

\begin{definition}\label{power_closed_below}
    $\powerClosedBelow{\nu}{\kappa}$ iff
        for every cardinal $\mu\in\kappa$ we have
            $\cardPower{\mu}{\nu}\in\kappa$.
\end{definition}

\begin{proposition}\label{gch_point_elim}
    Suppose $\gchPoint{\mu}$.
    Then $\mu$ is an infinite cardinal and
        $\cardPower{\cardTwo}{\mu}=\cardSucc{\mu}$.
\end{proposition}
\begin{proof}
    Follows by \obv{gch_point}.
\end{proof}

\begin{proposition}\label{gch_points_subseteq}
    $\gchPoints{\kappa}\subseteq\kappa$.
\end{proposition}
\begin{proof}
    Follows by \obv{gch_points,subseteq}.
\end{proof}

\begin{proposition}\label{gch_points_elim}
    Let $\mu\in\gchPoints{\kappa}$.
    Then $\mu\in\kappa$ and $\gchPoint{\mu}$.
\end{proposition}
\begin{proof}
    Follows by \obv{gch_points}.
\end{proof}

\subsection{The local power bound}

\begin{theorem}[Local power closure]\label{stationary_gch_points_imply_power_closed_below_cofinality}
    Let $\kappa$ be a singular cardinal.
    Let $\lambda=\cofinality{\kappa}$.
    Suppose $\gchPoints{\kappa}$ is stationary in $\kappa$.
    Then $\powerClosedBelow{\lambda}{\kappa}$.
\end{theorem}
\begin{proof}
    $\kappa$ is a limit ordinal
        by \obv{infinite_cardinal_is_limit_ordinal}.
    Hence $\kappa$ is an ordinal by \obv{limit_ordinal}.
    $\lambda$ is a cardinal by \obv{cofinality_is_cardinal}.
    Hence $\lambda$ is an ordinal by \auto{cardinal,initial_ordinal}.
    $\lambda\in\kappa$
        by \auto{cofinality_belongs_to_singular_cardinal}.
    $\ordinalUnbounded{\gchPoints{\kappa}}{\kappa}$
        by \obv{stationary_is_unbounded}.
    Show that for every cardinal $\mu\in\kappa$
        we have $\cardPower{\mu}{\lambda}\in\kappa$.
    \begin{subproof}
        Fix $\mu$.
        Assume $\mu\in\kappa$.
        Assume $\mu$ is a cardinal.
        Hence $\mu$ is an ordinal by \auto{cardinal,initial_ordinal}.
        Let $\tau=\mu\union\lambda$.
        $\mu\subseteq\lambda$ or $\lambda\subseteq\mu$
            by \obv{ordinal_subseteq_or_subseteq}.
        Hence $\tau=\lambda$ or $\tau=\mu$
            by \obv{union_absorb_subseteq_left,union_absorb_subseteq_right}.
        Hence $\tau\in\kappa$.
        Take $\theta\in\gchPoints{\kappa}$ such that $\tau\in\theta$
            by \obv{ordinal_unbounded}.
        Hence $\theta\in\kappa$ and $\gchPoint{\theta}$
            by \obv{gch_points_elim}.
        Hence $\theta$ is an infinite cardinal and
            $\cardPower{\cardTwo}{\theta}=\cardSucc{\theta}$
            by \obv{gch_point_elim}.
        Hence $\theta$ is an ordinal by \auto{cardinal,initial_ordinal}.
        $\mu\subseteq\tau$ and $\lambda\subseteq\tau$
            by \obv{union_upper_left,union_upper_right}.
        $\tau\subseteq\theta$ by \obv{ordinal_elem_implies_subseteq}.
        Hence $\mu\subseteq\theta$ and $\lambda\subseteq\theta$
            by \obv{subseteq_transitive}.
        Hence $\cardLe{\mu}{\theta}$ and $\cardLe{\lambda}{\theta}$
            by \obv{subseteq_cardinal_le}.
        Hence $\cardLe{\cardPower{\mu}{\lambda}}
            {\cardPower{\theta}{\lambda}}$
            by \obv{cardinal_power_monotone_base}.
        $\emptyset\in\theta$ by \obv{zero_one_in_infinite_cardinal}.
        Hence $\theta$ is inhabited.
        Hence $\cardLe{\cardPower{\theta}{\lambda}}
            {\cardPower{\theta}{\theta}}$
            by \obv{cardinal_power_monotone_exponent}.
        Hence $\cardLe{\cardPower{\mu}{\lambda}}
            {\cardPower{\theta}{\theta}}$
            by \obv{cardinal_le_transitive}.
        $\cardPower{\mu}{\lambda}$ and $\cardPower{\theta}{\theta}$
            are initial ordinals by \auto{cardinal_power_is_cardinal,cardinal}.
        \begin{align*}
            \cardPower{\mu}{\lambda}
                &\subseteq \cardPower{\theta}{\theta}
                    \explanation{by \obv{initial_ordinal_cardinal_le_iff_subseteq}}
                \\
                &= \cardPower{\cardTwo}{\theta}
                    \explanation{by \obv{infinite_cardinal_self_power}}
                \\
                &= \cardSucc{\theta}
        \end{align*}
        $\cardPower{\mu}{\lambda}$ is an ordinal by \obv{initial_ordinal}.
        $\cardSucc{\theta}$ is an ordinal
            by \auto{cardinal_successor_is_cardinal,cardinal,initial_ordinal}.
        $\cardSucc{\theta}\in\kappa$
            by \obv{infinite_cardinal_successor_belongs_to_singular_cardinal}.
        Hence $\cardPower{\mu}{\lambda}\in\cardSucc{\theta}$ or
            $\cardPower{\mu}{\lambda}=\cardSucc{\theta}$
            by \obv{ordinal_subseteq_implies_elem_or_eq}.
        Follows by \auto{ordinal_transitivity}.
    \end{subproof}
    Follows by \obv{power_closed_below}.
\end{proof}

\subsection{Counting subsets of the index set}

\begin{proposition}\label{power_closed_below_powerset}
    Suppose $\powerClosedBelow{\nu}{\kappa}$.
    Suppose $\cardTwo\in\kappa$.
    Then $\card{\pow{\nu}}\in\kappa$.
\end{proposition}
\begin{proof}
    $\cardTwo$ is a cardinal by \obv{cardinal_two_is_cardinal}.
    Hence $\cardPower{\cardTwo}{\nu}\in\kappa$
        by \obv{power_closed_below}.
    $\cardPower{\cardTwo}{\nu}=\card{\pow{\nu}}$
        by \obv{cardinal_power_two_is_powerset_cardinal}.
\end{proof}

\begin{theorem}\label{stationary_gch_points_bound_cofinality_powerset}
    Let $\kappa$ be a singular cardinal.
    Let $\lambda=\cofinality{\kappa}$.
    Suppose $\gchPoints{\kappa}$ is stationary in $\kappa$.
    Then $\card{\pow{\lambda}}\in\kappa$.
\end{theorem}
\begin{proof}
    $\cardOne\in\naturals$
        by \auto{cardinal_one,omega_empty,omega_successor}.
    Hence $\cardTwo\in\naturals$
        by \auto{cardinal_two_successor,omega_successor}.
    $\naturals\subseteq\kappa$
        by \obv{naturals_subseteq_infinite_cardinal}.
    Hence $\cardTwo\in\kappa$ by \obv{subseteq}.
    $\powerClosedBelow{\lambda}{\kappa}$
        by \obv{stationary_gch_points_imply_power_closed_below_cofinality}.
    Follows by \obv{power_closed_below_powerset}.
\end{proof}
